% M27, 13.09.2026: Einleitung aus dem konsolidierten Hauptteil.
% Overleaf-Ziel: Kapitel/01/01-Einleitung.tex (vollständige Datei ersetzen).
% Quellen und Fragenzuordnung: M27-Umsetzung-2026-09-13.md.
% Allgemeine Motivation und explorative Entstehung sind getrennt;
% kein vorab festgelegtes Gesamtdesign, kein Effizienzversprechen.

\chapter{Einleitung}
\label{chap:einleitung}

\section{Motivation und Problemstellung}
\label{sec:einleitung-motivation}

In den frühen Phasen eines Softwareprojekts wird geklärt, welches
Problem ein System lösen soll, welche Funktionen benötigt werden und
welche Randbedingungen seine Gestaltung bestimmen. Diese Informationen
geben der weiteren Entwicklung eine Richtung. Ihre Qualität ist deshalb
auch für spätere Arbeiten relevant. Kasauli et al. beschreiben in einer
Mehrfachfallstudie mit sieben Unternehmen der großskaligen agilen
Systementwicklung das Risiko umfangreicher Nacharbeit, wenn etwa
Leistungs- oder Gebrauchstauglichkeitsprobleme erst spät bearbeitet
werden. Zugleich arbeiten sie einen Zielkonflikt heraus: Fehlende
Anforderungsinformationen können die Entwicklung verzögern, während
eine übermäßig detaillierte Ausarbeitung ebenfalls Zeit beansprucht
\cite[S.~1, 19]{kasauli2021largescaleagilere}.
Für die vorliegende Untersuchung motiviert dies eine Informationsführung,
die frühe Klärung unterstützt und spätere Erkenntnisse aufnehmen kann.

Ein Gespräch zu Projektbeginn liefert dafür eine mögliche Quelle.
Es enthält Wünsche, Begründungen, offene Fragen und erste
Lösungsvorstellungen, die zunächst eingeordnet und miteinander in
Beziehung gesetzt werden müssen. Die anschließende Planung benötigt
unterschiedliche Ergebnisse: Anforderungen beschreiben benötigtes
Verhalten und Qualitäten, Architekturinformationen betreffen die technische Gestaltung,
und Arbeitspakete bereiten die Umsetzung vor. Werden Inhalte dabei
verdichtet oder umformuliert, stellt sich die Frage, auf welchen
Gesprächsaussagen die Ergebnisse beruhen. Kommen später Ergänzungen oder
Korrekturen hinzu, muss außerdem erkennbar bleiben, welche bestehenden
Inhalte sie betreffen und wie mit ihnen weitergearbeitet wird.

Große Sprachmodelle, im Folgenden \emph{Large Language Models} (LLMs),
eröffnen Möglichkeiten, solche sprachlich vorliegenden Informationen
zu verarbeiten. Die systematische Literaturübersicht von Cheng et al.
dokumentiert Anwendungen generativer KI insbesondere bei der Erhebung,
Analyse und Spezifikation von Anforderungen. Die fortlaufende
Anforderungsverwaltung wird in dem untersuchten Literaturbestand
seltener behandelt; zugleich werden Grenzen der Reproduzierbarkeit,
Verlässlichkeit und praktischen Erprobung deutlich
\cite[S.~141, 153, 165]{cheng2026generative}.
Ein geeigneter Untersuchungsgegenstand ist damit die Verbindung von
inhaltlicher Verarbeitung und weiterführender Bearbeitung im
Projektzusammenhang.

LLM-basierte KI-Agenten betten die Modellverarbeitung in
zielgerichtete Handlungen ein. Sie können beispielsweise Teilaufgaben
bearbeiten, Werkzeuge auswählen oder Rückmeldungen in weitere Schritte
einbeziehen. Mehrere spezialisierte Agenten lassen sich dazu über eine
Orchestrierung verbinden, also über eine Koordination ihrer Aufgaben
und Informationsflüsse. He et al. beschreiben entsprechende Ansätze
auch für die Bearbeitung mehrerer Aufgaben der Anforderungsentwicklung,
von der Erhebung über die Prüfung bis zur Spezifikation
\cite[S.~124:4, 124:6]{he2025llmbasedmultiagent}.
Für ein daraus aufgebautes System muss jedoch festgelegt werden,
welche Entscheidungen Agenten treffen dürfen, wie Ergebnisse geprüft
werden und wann sie in den verbindlichen Projektbestand eingehen.

Hier setzt die Arbeit an. Sie untersucht, wie agentische Verarbeitung
so in einen Gesamtprozess eingebunden werden kann, dass aus
Projektkommunikation fachliche Ergebnisse entstehen und deren Herkunft,
Beziehungen und Änderungen untersucht werden können. Der in
Abschnitt~\ref{sec:forschungsstand-positionierung} ausgeführte
Forschungsvergleich ordnet diesen Untersuchungsbedarf gegenüber
bestehenden Ansätzen ein. Im Mittelpunkt steht die konkrete Verbindung
von Quellenbezug, Projektzustand und kontrollierter Fortschreibung.

\section{Zielsetzung und Forschungsfragen}
\label{sec:einleitung-forschungsfragen}

Ziel ist die Entwicklung und Untersuchung eines Agentensystems, das
Projektinformationen in Anforderungen, Architekturinformationen und
Arbeitspakete überführt und bei weiteren Eingaben auf vorhandene
Ergebnisse zurückgreift. Der betrachtete Ausschnitt des
\emph{Software Development Lifecycle} (SDLC), also des Lebenszyklus
eines Softwaresystems, umfasst damit die Informationsarbeit für
Anforderungsanalyse, frühe Architekturplanung und die Vorbereitung der
Umsetzung. Microsoft Agent Framework (MAF) ist der gesetzte technische
Rahmen für die Realisierung.

Das Ausgangsszenario macht diese Aufgabe konkret: Ein
Kick-off-Transkript behandelt die Planung einer App für eine
Einrichtung für Menschen mit Beeinträchtigungen. Die App soll
Betreuenden unter anderem Wissen über die individuelle Kommunikation
betreuter Personen zugänglich machen. Diese zu planende Anwendung ist
vom Forschungsartefakt zu unterscheiden: Untersucht wird das
Agentensystem, das ihre Planungsinformationen verarbeitet.
Abschnitt~\ref{sec:ausgangslage-entwicklungsbefunde} erläutert das
Ausgangsmaterial und den Projektkontext.

Die Hauptforschungsfrage lautet:
\begin{quote}
\emph{Wie können KI-Agenten mittels Microsoft Agent Framework in den
frühen Phasen des Software Development Lifecycle eingesetzt werden?}
\end{quote}

Für die Bearbeitung werden drei Teilfragen unterschieden:
\begin{description}
  \item[TF1:] Welche Gestaltungsprobleme und Designanforderungen ergeben
  sich aus dem angestrebten Einsatz und der explorativen Entwicklung
  des Agentensystems?
  \item[TF2:] Wie lassen sich diese Anforderungen durch die Verteilung
  von Aufgaben und Entscheidungen zwischen Agenten, Ablaufsteuerung
  und menschlicher Mitwirkung adressieren und mit MAF realisieren?
  \item[TF3:] Welche Eigenschaften, Zielkonflikte und Grenzen zeigen
  die ausgewählten Untersuchungen der Inhaltsverarbeitung, der
  nachvollziehbaren Fortschreibung und der technischen
  Kontrollmechanismen?
\end{description}

Die Teilfragen gliedern den Erkenntnisweg der Arbeit. TF1 wird anhand
der ausgewählten Entwicklungsbefunde und der Anforderungsherleitung
bearbeitet, TF2 anhand der Architektur und ihrer Realisierung.
TF3 richtet sich auf die empirischen und technischen Nachweise der
Evaluation. Tabelle~\ref{t:methodischer-zugang-teilfragen} konkretisiert
die jeweiligen methodischen Zugänge;
Kapitel~\ref{chap:diskussion-fazit} führt die Antworten zusammen.

\section{Vorgehen, Beitrag und Abgrenzung}
\label{sec:einleitung-vorgehen-beitrag}

Die Arbeit verfolgt einen explorativen, artefaktzentrierten Ansatz mit
Elementen der Design Science Research, also einer Forschung durch die
Gestaltung und Untersuchung eines Artefakts. Zu Beginn stand ein
einfaches technisches Ziel: Ein Transkript sollte durch KI-Agenten mit
MAF verarbeitet werden, sodass daraus ein GitHub-Issue als operative
Arbeitsaufgabe entsteht. Die spätere Erweiterung um Rückkopplung und
Wiederverwendung gab Anlass, auch die weitere Bearbeitung vorhandener
Ergebnisse zu untersuchen. Die konkrete Aufgabenverteilung,
Informationsstruktur und Ablaufsteuerung waren dabei zunächst offen.
Sie entstanden schrittweise aus Prototypen, beobachteten Problemen und
Gestaltungsentscheidungen. Kapitel~\ref{chap:forschungsmethodik}
erläutert die wissenschaftliche Einordnung dieses Vorgehens.

Die Rekonstruktion der Entwicklung konzentriert sich auf Beobachtungen,
die tragende Anforderungen und Architekturentscheidungen erklären.
Diese formative Untersuchung dient der Begründung des Entwurfs.
Die anschließende Evaluation charakterisiert ausgewählte Eigenschaften
des realisierten Systems. Hierfür werden Inhaltsausgaben mit
Referenzaussagen verglichen, gespeicherte Zwischenstände und
Fortschreibungsfälle ausgewertet sowie technische Kontrollmechanismen
an definierten Prüffällen untersucht. Die jeweiligen System- und
Auswertungsstände werden getrennt ausgewiesen. Entwicklungsbeobachtungen
und technische Implementierung bilden somit jeweils andere Belege als
eine vergleichende Inhaltsmessung.

Der Beitrag liegt in der begründeten Integration dieser Aufgaben in
einem ausführbaren Forschungsartefakt und in der Untersuchung seiner
Möglichkeiten und Grenzen. Die Architektur verbindet semantische
Modellarbeit, etwa die Interpretation und fachliche Einordnung von
Aussagen, mit expliziten Workflows. Diese legen die Verarbeitungsschritte
und ihre Übergänge fest. Regelprüfungen sichern formal prüfbare
Bedingungen; menschliche Entscheidungen bestimmen die fachliche
Übernahme in den Fortschreibungspfaden. Strukturierte Zwischenstände
und ein über einzelne Ausführungen hinaus geführter Projektbestand
machen Herkunft und Änderungen zu Gegenständen der Prüfung.
Die MAF-Realisierung zeigt, wie die Frameworkmechanismen und die
projektseitig entwickelten Informations- und Kontrollstrukturen
zusammenwirken.

Die Untersuchung konzentriert sich auf einen begrenzten Projekt- und
Bedienkontext mit ausgewählten Quellen, Verarbeitungsläufen und
Prüffällen. Die Implementierung der geplanten App und die nachfolgenden
SDLC-Phasen gehören nicht zum Untersuchungsumfang. Die Evaluation
liefert zudem keinen Nachweis einer allgemeinen Überlegenheit gegenüber
anderen Agentenarchitekturen oder Frameworks. Tokenangaben
charakterisieren den technischen Aufwand der untersuchten Ausführungen;
eine Verringerung der gesamten Projektkosten oder des menschlichen
Prüfaufwands wird damit nicht nachgewiesen. Nachvollziehbarkeit,
Inhaltserhaltung und technische Kontrolle werden als unterschiedliche
Eigenschaften bewertet.

\section{Aufbau der Arbeit}
\label{sec:einleitung-aufbau}

Kapitel~\ref{chap:grundlagen-forschungsstand} erläutert die fachlichen
und technischen Grundlagen und ordnet verwandte Forschungsansätze ein.
Kapitel~\ref{chap:forschungsmethodik} beschreibt, wie die explorative
Entwicklung und die Evaluation zur Beantwortung der Forschungsfragen
beitragen.

Die Kapitel~\ref{chap:explorative-problemklaerung} bis
\ref{chap:maf-realisierung} bilden den zusammenhängenden
Gestaltungs- und Realisierungsbogen:
Kapitel~\ref{chap:explorative-problemklaerung} leitet aus Nutzungsziel
und Entwicklungsbefunden die Designanforderungen her.
Kapitel~\ref{chap:loesungsarchitektur} erklärt die daraus entwickelte
Architektur und ihre Funktionsweise.
Kapitel~\ref{chap:maf-realisierung} vertieft die dafür relevanten
Realisierungsentscheidungen mit MAF.

Kapitel~\ref{chap:evaluation} berichtet die ausgewählten Untersuchungen
und ihre Ergebnisse. Kapitel~\ref{chap:diskussion-fazit} führt diese mit
der Gestaltungsbegründung zusammen, beantwortet die Forschungsfragen
und diskutiert Beitrag, Grenzen und weiterführende Untersuchungen.
